%% notes-tex/028-gene-essentiality/sections/7-closing.tex
%% [[experiments.028-gene-essentiality]]
%% The rounds run after the first summary was drawn: each fixes its read in the pack file
%% before its cells finish, and each table is generated from the cells' result files.

\section{The closing rounds: repeated splits, the rules, and what each simplification costs\secstatus{todo}}
\label{sec:closing}

The released held-out set has 31 essential genes, so every conclusion drawn on it alone
is a conclusion about 31 genes. The rounds in this section reread the finalists and the
module's controls on ten other held-out fifths of the labeled genes, add the
gene-protein-reaction rules beside the solve, and price each way of making the module
cheaper or of fitting it inside one experiment per record.

\subsection{Round 75: the finalists on ten repeated outer splits\secstatus{todo}}
\label{sec:round75}

The released held-out set has 31 essential genes, the finalists' intervals on it overlap,
and the average's 0.954 is the maximum of eleven averages chosen on those genes. Round 75
(jobs 2547 to 2567; Table~\ref{tab:resplit}) reads the same models on ten other held-out
sets, each a class-stratified fifth of the released labeled genes with validation a
stratified fifth of the rest (\texttt{build\_essentiality\_splits.py --resplit}), three
seeds per model per split. The models, the equal-weight average and the statistics were
committed before any cell finished (\texttt{resplit\_readout.py}).

%% notes-tex/028-gene-essentiality/tables/resplit_readout.tex
%% GENERATED by experiments/028-gene-essentiality/scripts/notes_tex_tables.py
%% SOURCE: results/resplit_readout.json
\begin{table}[H]
  \centering
  \footnotesize
  \caption[Round 75: ten repeated outer splits]{Round 75, the finalists on ten repeated outer splits. Each test block is a class-stratified fifth of the released labeled genes (198 genes, 32 essential), validation a stratified fifth of the rest. Left: the AUROC of the three-seed family score (standardized logits averaged over seeds) of the transformer with the layer-free solve, of graph convolution on GO at 400 epochs, and of their equal-weight average, fixed before any cell finished; best of the three in bold. Right: the seed-mean AUROC per cell of the layer-free cell and its two controls on the same backbone. Below: the paired differences across the ten splits. The test blocks overlap, so the differences are not independent and no $p$-value is given. Written by \texttt{resplit\_readout.py}.}
  \label{tab:resplit}
  \begin{tabular}{@{}lrrrrrr@{}}
    \hline
    split & layer-free & GO conv. & average & layer-free, cells & FBA ratios, cells & no metabolism, cells \\
    \hline
    0 & 0.864 & 0.884 & \textbf{0.901} & 0.863 & 0.864 & 0.840 \\
    1 & 0.844 & 0.884 & \textbf{0.899} & 0.843 & 0.803 & 0.794 \\
    2 & 0.857 & 0.872 & \textbf{0.893} & 0.854 & 0.844 & 0.854 \\
    3 & 0.850 & 0.857 & \textbf{0.895} & 0.849 & 0.835 & 0.774 \\
    4 & 0.826 & \textbf{0.868} & 0.868 & 0.824 & 0.843 & 0.773 \\
    5 & 0.857 & 0.868 & \textbf{0.892} & 0.849 & 0.830 & 0.741 \\
    6 & 0.857 & 0.889 & \textbf{0.911} & 0.842 & 0.835 & 0.814 \\
    7 & 0.786 & 0.838 & \textbf{0.844} & 0.783 & 0.737 & 0.729 \\
    8 & 0.904 & 0.872 & \textbf{0.936} & 0.899 & 0.878 & 0.826 \\
    9 & 0.831 & 0.860 & \textbf{0.918} & 0.824 & 0.853 & 0.821 \\
    \hline
    mean & 0.848 & 0.869 & 0.896 & 0.843 & 0.832 & 0.796 \\
    \hline
    \multicolumn{4}{@{}l}{average minus layer-free} & \multicolumn{3}{l}{$+0.048$ (sd 0.017), positive on 10 of 10} \\
    \multicolumn{4}{@{}l}{average minus GO convolution} & \multicolumn{3}{l}{$+0.026$ (sd 0.021), positive on 9 of 10} \\
    \multicolumn{4}{@{}l}{layer-free minus GO convolution} & \multicolumn{3}{l}{$-0.022$ (sd 0.024), positive on 1 of 10} \\
    \multicolumn{4}{@{}l}{average minus the better single model} & \multicolumn{3}{l}{$+0.023$ (sd 0.017), positive on 9 of 10} \\
    \multicolumn{4}{@{}l}{layer-free minus FBA ratio affine (seed means)} & \multicolumn{3}{l}{$+0.011$ (sd 0.023), positive on 7 of 10} \\
    \multicolumn{4}{@{}l}{layer-free minus no metabolism (seed means)} & \multicolumn{3}{l}{$+0.047$ (sd 0.034), positive on 10 of 10} \\
    \hline
  \end{tabular}
\end{table}


The average is the best method and the margin is its own, not a selection. It reads 0.896
as the mean of ten splits against 0.869 for graph convolution on GO and 0.848 for the
transformer with the layer-free solve; it is above the transformer on ten of ten
($+0.048$, sd 0.017), above the graph model on nine ($+0.026$, sd 0.021), and above the
better of the two on nine ($+0.023$, sd 0.017). The two single models are not the tie the
released genes showed: graph convolution on GO is above the transformer on nine of ten
($-0.022$, sd 0.024 for the transformer minus the graph model).

With the single-experiment form of Section~\ref{sec:round84-98} as the transformer family
(the YPD ratio, the flux vector and the rule counts, six seeds per split, against graph
convolution on GO at six), the same rule reads 0.907 for the average against 0.871 for
the graph model and 0.869 for the transformer. The average is above the better of the two
on ten of ten splits ($+0.029$, sd 0.013), and the two single models tie ($-0.003$, the
transformer higher on four of ten). On the released genes its five-fold ensemble reads 0.929 alone and 0.953 [0.92, 0.98]
averaged with the graph family's, against 0.955 for the four-media recipe; that is one
more post hoc average on those genes and it is not the maximum. The form that goal 2 keeps is therefore also the
better partner for goal 1.

Every model reads lower on these fifths than on the released held-out genes, the
transformer's cells 0.843 against 0.908 and the graph model's 0.868 against 0.905 as seed
means, so the
released held-out set ranks more easily than a random fifth of the labeled genes. What
makes it so is not measured.

The module's contribution splits in two on the same backbone. The solve as an input is
worth $+0.047$ over no metabolism, on ten of ten splits (sd 0.034). Within that, the
deletion's flux vector over FBA's four growth ratios alone is $+0.011$ (sd 0.023, seven
of ten), where the released split read $+0.018$ (0.908 against 0.890): the sign agrees
and the difference does not separate from zero. Doubling the seeds does not separate it
either (round 78, jobs 2586 to 2595): at six seeds per arm per split the difference is
$+0.010$ (sd 0.025, eight of ten), the same mean with the same spread, so the power was
not what the three-seed read lacked. So what the ten splits support is that the solve
helps and that the four ratios carry most of it, while the 4,131 flux columns add
little beyond them; six counts of gene-protein-reaction redundancy, on the same splits,
add about twice as much ($+0.019$, Section~\ref{sec:round77}).

\subsection{Round 76: the physical graph's degree as an input to the transformer\secstatus{todo}}
\label{sec:round76}

Off the metabolic genes the raw degree in the physical graph ranks the 4,674 genome-wide
genes at 0.740 with nothing trained, and the transformer with the layer-free solve read
0.570. Round 76 (jobs 2568 to 2573, six seeds per arm; Table~\ref{tab:genome-heldout})
gives the transformer that degree the cheapest way, as two more affine columns beside
the flux vector, the gene's degree and clustering coefficient in the physical graph,
over a store that covers all 6,607 genes (\texttt{degree\_affine\_store.py}; a gene
with no deletion vector reads the wild-type vector instead of the all-zero row the
earlier store gave it).

The control changes nothing, as a constant row should: with wild-type rows and no degree
the six cells read the earlier arm's numbers to four decimals (mean 0.566, pooled 0.570).
With the two columns the cells read 0.564 to 0.656 (mean 0.602, sd 0.030; pooled 0.593
[0.57, 0.61]), above the control on six of six seeds ($+0.036$), and 0.16 under the
graph convolution's 0.755. On the released split the same store reads 0.906 (sd 0.009)
against 0.908, so the columns cost nothing there. The degree enters, and an affine
fitted on 629 metabolic genes takes a fraction of it: among the 787 released training
and validation genes degree ranks essentiality at 0.59, so the slope the affine can learn from them is small. Hypothesis,
untested: the gap to the convolution is the training set and not the path, and a
convolution trained on the same metabolic genes reaches 0.755 because its two layers
read the neighbors' degrees as well as the gene's own.

\subsection{Round 77: the gene-protein-reaction rules as affine columns\secstatus{todo}}
\label{sec:round77}

Section~\ref{sec:hard-genes} named what the solve cannot mark: the 19 essential held-out
genes whose knockout growth ratio is 1.0, four of them covered by an alternative gene set
on every reaction and fifteen holding a reaction of their own that the biomass objective
does not need. The flux vector is close to the wild type's on all nineteen either way, so
round 77 (jobs 2574 to 2585; Table~\ref{tab:gpr-readout}) appends what the vector does not
carry: six counts read off yeast-GEM's gene-protein-reaction rules in disjunctive normal
form (\texttt{isozyme\_affine\_store.py}), the gene's reactions, how many its deletion
alone disables, how many stay enabled, the minimum and mean number of alternative gene
sets, and its largest complex. All are non-negative, so the affine's $\log(1 + |v|)$
keeps their order and the columns cost six weights.

They help, and by the smaller of the two mechanisms the section named. On the released
split six seeds read 0.918 against 0.908 without them, above on six of six; on the ten
repeated outer splits the seed-mean per split is 0.862 against 0.843, above on eight of
ten ($+0.019$, sd 0.020). Over the 19 genes the mean misrank falls from 0.134 to 0.120.

The split inside those 19 is where the columns say what they are doing. The fifteen genes
holding a reaction of their own improve on six seeds of six, from 0.124 to 0.089, while
the four an alternative gene set covers get worse on six of six, from 0.171 to 0.234. The
columns state what the deletion's flux vector cannot, that a gene holds a reaction nothing
else catalyzes, because the program routes around those reactions without needing them;
and where the rules say a gene is covered the columns push the cell further toward calling
it dispensable, which is the wrong call on exactly those four. That pattern suggested the
count of sole reactions was the one useful column; round 84 tested it and the six columns
are redundant among themselves (Section~\ref{sec:round84-98}). The
hypothesis Section~\ref{sec:hard-genes} raised, that a paralog-aware input would move the
covered genes, is refuted on this construction: it moves the others.

%% notes-tex/028-gene-essentiality/tables/gpr_readout.tex
%% GENERATED by experiments/028-gene-essentiality/scripts/notes_tex_tables.py
%% SOURCE: results/gpr_readout.json
\begin{table}[H]
  \centering
  \footnotesize
  \caption[Round 77: the gene-rule redundancy columns]{Round 77, six columns of gene-protein-reaction redundancy appended to the deletion's flux vector in the affine (\texttt{isozyme\_affine\_store.py}: reactions, sole reactions, covered reactions, minimum and mean alternative gene sets, co-subunits), against the same cell without them. Rows one and two are held-out AUROC, six seeds on the released split and ten split means of three seeds each; rows three to five are the mean misrank of the essential held-out genes whose FBA knockout growth ratio is 1.0, for which a LOWER value is better and the sign count therefore counts the seeds made worse. Better of the two in bold. Written by \texttt{gpr\_readout.py}.}
  \label{tab:gpr-readout}
  \begin{tabular}{@{}p{62mm}rrrrl@{}}
    \hline
    quantity & with & without & difference & sd & positive \\
    \hline
    released split, held-out AUROC & \textbf{0.918} & 0.908 & +0.0102 & 0.0053 & 6 of 6 \\
    ten repeated outer splits, AUROC & \textbf{0.862} & 0.843 & +0.0192 & 0.0198 & 8 of 10 \\
    misrank, the 19 essential genes FBA calls viable & \textbf{0.120} & 0.134 & -0.0139 & 0.0111 & 1 of 6 \\
    of them, the 15 holding a reaction of their own & \textbf{0.089} & 0.124 & -0.0343 & 0.0133 & 0 of 6 \\
    of them, the 4 an alternative gene set covers & 0.234 & \textbf{0.171} & +0.0628 & 0.0288 & 6 of 6 \\
    \hline
  \end{tabular}
\end{table}


\subsection{Round 79: the class token through the perturbation operator\secstatus{todo}}
\label{sec:round79}

Every pooled cell in this document reads the genotype through a mean over the perturbed
gene tokens, because the class token was sliced off BEFORE the perturbation operator and
is therefore the encoding of the unperturbed graph, the same vector for every strain. A
flag that sends it through the operator as one more query over the same deleted-gene keys
existed on another branch, unlanded and absent from this worktree until it was ported on
2026-09-20; it adds no parameters and leaves the gene rows bit-identical. Round 79 (jobs
2597 to 2600; Table~\ref{tab:cls-readout}) is the paired read on the released split.

The token becomes strain-specific and the ranking gets slightly worse. Its across-strain
spread ends at 0.52 against exactly 0 without the flag, so the mechanism is doing what it
claims. On the seeds of round 79 alone neither arm separated from zero; round 81 (jobs
2603 to 2605) completed the pooled arm's pairing and took the layer-free arm to nine
pairs, and at fifteen pairs the cost is established: the layer-free cell reads 0.905
against 0.911, a paired difference of $-0.0054$ (95 percent interval $-0.0100$ to
$-0.0009$, two of nine above), the pooled cell 0.855 against 0.858 ($-0.0028$, one of
six above), and the two together $-0.0044$ (interval $-0.0075$ to $-0.0012$, twelve of
fifteen below). It is a small cost, about a third of this cell's seed spread, and it is
consistent.

The construction explains it. A strain here deletes ONE gene, so the class token's query
attends over exactly one key and the perturbed token is a function of the wild-type token
and that gene's token, which the perturbed-gene pool already carries. The flag adds no
information on a single deletion and widens the head's input, which is the shape of a
small consistent cost rather than a gain. It has something to add only where a genotype
perturbs several genes and the pool's mean over them discards their arrangement, which is
the setting it was written for and not this one. Here it stays off.

%% notes-tex/028-gene-essentiality/tables/cls_readout.tex
%% GENERATED by experiments/028-gene-essentiality/scripts/notes_tex_tables.py
%% SOURCE: results/cls_readout.json
\begin{table}[H]
  \centering
  \footnotesize
  \caption[Round 79: the perturbed class token]{Round 79, the class token sent through the perturbation operator as query row 0 over the same deleted-gene keys (\texttt{perturb\_cls}, ported from the 025 branch), against the identical cell without it, paired by seed on the released split. No parameters are added and the gene rows are unchanged. The last column is the across-strain spread of the class token at the final epoch, which is exactly 0 without the flag because the token is then the encoding of the unperturbed graph. Written by \texttt{cls\_readout.py}.}
  \label{tab:cls-readout}
  \begin{tabular}{@{}p{52mm}rrrrrlr@{}}
    \hline
    arm & seeds & with & without & difference & sd & positive & token sd \\
    \hline
    layer-free solve vector in the affine & 9 & 0.905 & \textbf{0.911} & -0.0054 & 0.0059 & 2 of 9 & 0.52 \\
    pooled cell, no solve and no affine & 6 & 0.855 & \textbf{0.858} & -0.0028 & 0.0054 & 1 of 6 & 0.60 \\
    \hline
  \end{tabular}
\end{table}


\subsection{Round 80: the rule counts under the reported protocol\secstatus{todo}}
\label{sec:round80}

Round 77 read the gene-rule columns on single cells. Round 80 (jobs 2601 and 2602) runs
the same cell on the five folds, so the document's best module form has a number on the
protocol the reported 0.893 was computed under, and so the two-family average can be
formed from it.

It is the highest single model in the document on both counts. The five fold models read
0.9335 [0.890, 0.968] as an ensemble (fold mean 0.9133), against the layer-free form's
0.920 and graph convolution on GO's 0.917, and 0.930 as a standardized family. Averaged
with the GO convolution's fold models it reads 0.9550 [0.922, 0.982]
(Table~\ref{tab:family-ensembles-folds}), a little above the 0.9538 of the cached-solve
average and the 0.9532 of the layer-free one. That number is now the maximum of thirteen
post hoc averages rather than eleven, and carries the selection that implies; the three
best are within 0.002 of one another and their intervals are nearly identical, so what
the folds establish is that the two families average above either alone, not an ordering
among the three module forms.

\subsection{Round 82: what each half of the module is worth\secstatus{todo}}
\label{sec:round82}

Round 77 added six gene-rule counts beside the 4,131-column flux vector and round 78 found
the vector itself worth little over the four growth ratios, so the two halves have never
been measured against the same baseline. Round 82 (jobs 2606 to 2609) runs the counts with
NO flux vector, six seeds and the five folds, completing the factorial
(Table~\ref{tab:module-factorial}).

Both halves help and the solve is the larger of the two. Against the four growth ratios
alone the counts are worth $+0.012$ and the vector $+0.024$, paired on the three seeds all
four arms share, and together they reach 0.918. The fold ensembles order the same way:
0.909 for the counts alone, 0.920 for the vector alone, 0.934 for both. The counts are
worth more once the vector is already there than the vector is worth once the counts are
($+0.020$ against $+0.008$), which is what two partly overlapping signals look like.

This corrects a comparison made in passing after round 78. The counts' $+0.019$ there was
measured ON TOP of the vector while the vector's $+0.010$ was measured on top of the
ratios, so the two are different conditionals and cannot be read as a ranking. Against
one baseline the solve is worth about twice the counts on the released split, and the ten
repeated splits put them within each other's spread. What survives both readings is that
each half is worth about a hundredth to two hundredths of AUROC, that they are
complementary, and that the pair is the best form of the module.

Neither reading changes the practical conclusion. The counts cost no solver, no medium
and no per-genotype computation, and alone they carry half of what the module contributes
here; the solve carries the other half and needs one linear program per genotype and
medium, run once offline.

%% notes-tex/028-gene-essentiality/tables/module_factorial.tex
%% GENERATED by experiments/028-gene-essentiality/scripts/notes_tex_tables.py
%% SOURCE: results/ (the cells named in the rows)
\begin{table}[H]
  \centering
  \footnotesize
  \caption[Round 82: the solve and the rules, each alone]{Round 82, what each half of the module is worth on the same backbone, the pathway multi-hot with ProtT5 and the four FBA growth ratios on the logit. The solved flux vector is 4,131 columns from one linear program per genotype; the gene-rule counts are six columns read off yeast-GEM's rules with no solver. Columns: the number of seeds, their mean held-out AUROC and its spread, and the five-fold ensemble where one was run. Best of each column in bold. Written by \texttt{notes\_tex\_tables.py} from the cells' own result files.}
  \label{tab:module-factorial}
  \begin{tabular}{@{}lrrrr@{}}
    \hline
    affine columns & seeds & cell mean & sd & fold ensemble \\
    \hline
    four growth ratios only & 3 & 0.890 & 0.005 & \multicolumn{1}{c}{--} \\
    + gene-rule counts, no solve & 6 & 0.902 & 0.005 & 0.909 \\
    + solved flux vector, no counts & 9 & 0.911 & 0.007 & 0.920 \\
    + both & 9 & \textbf{0.920} & 0.006 & \textbf{0.934} \\
    \hline
  \end{tabular}
\end{table}


\subsection{Round 83: every column inside one experiment\secstatus{todo}}
\label{sec:round83}

A record in the database is one genotype in one environment, and the four FBA growth
ratios are four different media, so a cell carrying all four reads four experiments at
once and cannot be written as one record. The other metabolic columns are already
admissible: the solved flux vector is a function of genotype and medium and this store is
YPD, and the gene-rule counts carry no environment at all. Round 83 (jobs 2610 to 2615,
six seeds per arm) keeps the YPD ratio, drops the other three, and reads the same three
recipes (Table~\ref{tab:single-experiment}).

The restriction is free wherever the solve is present. The full recipe reads 0.916 against
0.918 at the round's six seeds (the table's row carries the nine seeds of round 107,
0.912 against 0.920, Section~\ref{sec:round84-98}), the vector without the counts 0.912 against 0.911, and only the arm with nothing
else loses anything: the ratio alone reads 0.877 against 0.890. So about a hundredth of
what the four ratios contributed was the model reading four environments, and the deletion's
own flux vector supplies it from within one.

That settles the question the endpoint could not answer on its own. Nothing in the module
as it now stands needs more than one experiment per record, so the form carries to a
database where the environment is part of the key rather than part of the feature vector.

%% notes-tex/028-gene-essentiality/tables/single_experiment.tex
%% GENERATED by experiments/028-gene-essentiality/scripts/notes_tex_tables.py
%% SOURCE: results/ (the cells named in the rows)
\begin{table}[H]
  \centering
  \footnotesize
  \caption[Round 83: the single-experiment ablation]{Round 83, every affine column restricted to what ONE experiment defines. The four FBA growth ratios are four different media, so a cell carrying all four reads four experiments at once, which a genotype-by-environment record cannot; the solved flux vector is already one medium (YPD) and the gene-rule counts carry no environment at all. The right block keeps the YPD ratio and drops the other three. Seed means of the held-out AUROC, better of each row in bold. Written by \texttt{notes\_tex\_tables.py} from the cells' own result files.}
  \label{tab:single-experiment}
  \begin{tabular}{@{}lrrrrr@{}}
    \hline
    affine columns & seeds & four media & seeds & YPD only & change \\
    \hline
    growth ratio alone & 3 & \textbf{0.890} & 6 & 0.877 & -0.012 \\
    + solved flux vector & 9 & 0.911 & 6 & \textbf{0.912} & +0.001 \\
    + vector and rule counts & 6 & \textbf{0.918} & 6 & 0.916 & -0.002 \\
    \hline
  \end{tabular}
\end{table}


\subsection{Round 108: DCell's ontology hierarchy over these splits\secstatus{todo}}
\label{sec:round108}

Graph convolution on the Gene Ontology multi-hot is the strongest single encoding here,
which raises the question of whether the ontology's own structure reads it better than
message passing does. DCell answers that structure directly: one small subsystem per GO
term, wired child to parent along the ontology, each term's state a function of its
children's states and of its own gene input, with a linear head on every term at weight
0.3 beside the root head.

As published it cannot be run on this endpoint. DCell gives every gene its own weight
into every term it belongs to, and a held-out gene's deletion never appears in training,
so those weights keep their initialization; measured, that arm reads 0.578 over six
seeds and 0.651 as a fold ensemble, and it loses to the graph on ten of ten repeated
splits (Table~\ref{tab:dcell-style}). Replacing the per-gene weights with one shared
weight on whether the deleted gene carries the term makes the same hierarchy able to rank
unseen genes, and it learns: 0.825 over six seeds, 0.848 as a fold ensemble. The
auxiliary per-term losses are worth $+0.026$ on the released split.

It still loses to message passing on the same annotations, by 0.067 over the ten repeated
splits, below on ten of ten. The hierarchy was sized on validation over ten
configurations, and the width the data supports is 8, well under the published rule of
$\max(20, 0.3 n)$ genes per term, whose uncapped form would give a widest subsystem of
1,651 units and 74.9 million effective parameters against 629 training genes; capped at
32 it is the worst configuration tried. So the ontology's grouping is not what carries
the GO signal on this endpoint, and the same annotations read through the gene graphs
carry it further.

\paragraph{The ontology of the trigenic DCell runs.} The DCell runs of experiments 005
and 006 read a different Gene Ontology from the 2,362-term multi-hot above: SGD's
annotations with those inferred from genetic interaction removed, redundant terms
collapsed, and terms kept at four or more genes, which leaves 2,655 subsystems. Built
here as a second store (\texttt{go006\_store.py}, jobs 2688 and 2689), it keeps 2,651
terms as columns, of which 1,421 are among the current 2,362; a labeled gene lights
34.5 of its columns on average against 46.5 in the current store, and the hierarchy is
flatter, 14 strata with a widest fan-in of 106 children against 15 and 24. Graph
convolution does not care which: on the ten repeated splits at three seeds it reads
$-0.002$ on the trigenic runs' ontology (higher on four of ten) and 0.905 against 0.909
on the released split. The hierarchy does: the shared-input DCell reads 0.714 against
0.841 on the released split over the same three seeds and $-0.020$ over the ten splits
(higher on two of ten, sd 0.041). Those reads used the configuration tuned on the current store; tuned on the trigenic
store over the same twelve configurations (job 2694), validation prefers width 8 with
normalization before the nonlinearity (0.876 against 0.873 for the current store's best),
and at that configuration the released split reads 0.709 over the same three seeds (job
2700), no better than the first read; the ten-split read at it is job 2701. The current
store is therefore the better of the two for the hierarchy, and neither ontology brings
it to the graph's level.

%% notes-tex/028-gene-essentiality/tables/dcell_style.tex
%% GENERATED by experiments/028-gene-essentiality/scripts/notes_tex_tables.py
%% SOURCE: results/dcell_style_*.json
\begin{table}[H]
  \centering
  \footnotesize
  \caption[Round 108: DCell over these splits]{Round 108, DCell's ontology hierarchy on this endpoint. Every term of the Gene Ontology gets a small subsystem whose state is a function of its children's states and its own gene input, wired along the same 2,362-term multi-hot the graph baselines read, with DCell's auxiliary per-term heads at weight 0.3. The published input gives each gene its own weight into each term, which a held-out gene's deletion never trains, so the first two rows replace it with one shared weight on whether the deleted gene carries the term. Columns: seeds on the released split, their mean held-out AUROC and its spread, the five-fold ensemble on the same 195 genes, and the paired difference against graph convolution on the same annotations over ten repeated outer splits with the count of splits the row wins. Test blocks overlap across splits, so no p-value. Width and normalization were chosen on validation over ten configurations. Written by \texttt{notes\_tex\_tables.py} from \texttt{dcell\_style\_baseline.py}.}
  \label{tab:dcell-style}
  \begin{tabular}{@{}lrrrrl@{}}
    \hline
    model & seeds & cell mean & sd & fold ensemble & against the graph, ten splits \\
    \hline
    DCell hierarchy, GO membership on each term & 6 & 0.825 & 0.026 & \textbf{0.848} & $-0.067$ (0.037), 0/10 \\
    the same, the module as an affine on the root & 3 & \textbf{0.846} & 0.017 & \multicolumn{1}{c}{--} & $-0.032$ (0.034), 2/10 \\
    the same, no auxiliary term losses & 6 & 0.799 & 0.037 & 0.844 & \multicolumn{1}{c}{--} \\
    DCell's published per-gene input & 6 & 0.578 & 0.057 & 0.651 & $-0.294$ (0.034), 0/10 \\
    \hline
    graph convolution on the same GO & 6 & 0.905 & 0.005 & 0.918 & 0.868 over ten splits \\
    \hline
  \end{tabular}
\end{table}


%% notes-tex/028-gene-essentiality/tables/go006_comparison.tex
%% GENERATED by experiments/028-gene-essentiality/scripts/notes_tex_tables.py
%% SOURCE: results/go006_comparison.json
\begin{table}[H]
  \centering
  \footnotesize
  \setlength{\tabcolsep}{2.5pt}
  \caption[The trigenic DCell runs' ontology against the current store]{The two GO-reading models on the ontology the trigenic DCell runs of experiments 005 and 006 used (2,651 columns: SGD annotations without those inferred from genetic interaction, redundant terms collapsed, four or more genes per term) and on the current store (2,362 columns: propagated annotations, ten or more genes per term). Three seeds per cell throughout, the current store's released value over the same three seeds. Columns: the released split, the mean over the ten repeated outer splits, the paired difference of the trigenic runs' ontology minus the current one with its spread over splits, and the splits it is higher on. Written by \texttt{notes\_tex\_tables.py} from \texttt{go006\_compare.py}.}
  \label{tab:go006-comparison}
  \begin{tabular}{@{}>{\raggedright\hangindent=1em\hangafter=1\arraybackslash}p{51.7mm}rrrrlr@{}}
    \hline
    model & released, trigenic & released, current & ten splits, trigenic & ten splits, current & difference & higher \\
    \hline
    graph convolution on GO & 0.905 & 0.909 & 0.867 & 0.868 & $-0.002$ (0.005) & 4/10 \\
    graph convolution on GO and the four ratios & 0.902 & 0.909 & 0.868 & 0.871 & $-0.004$ (0.011) & 2/10 \\
    graph convolution on ProtT5, GO and the ratios & 0.898 & 0.901 & 0.860 & 0.864 & $-0.003$ (0.004) & 2/10 \\
    DCell hierarchy, shared input & 0.714 & 0.841 & 0.781 & 0.801 & $-0.020$ (0.041) & 2/10 \\
    \hline
  \end{tabular}
\end{table}


\subsection{Rounds 109 and 110: the module on a backbone with no pathway annotation\secstatus{todo}}
\label{sec:round109}

Every transformer cell from round 75 onward reads the curated pathway multi-hot beside
ProtT5, and a dataset outside yeast may carry no such annotation. Rounds 109 and 110 (jobs
2685, 2686 and 2692) drop it: ProtT5 alone with the single-experiment module, the YPD
growth ratio, the YPD flux vector and the six gene-rule counts as the affine, against
ProtT5 alone with nothing on the logit. On the released split at six seeds the module
takes the cell from 0.771 to 0.883, $+0.112$ paired by seed and higher on six of six,
the largest gain of any pair in this document; the same pair on the pathway backbone reads
$+0.058$ (0.858 to 0.916). With the module the annotation-free backbone sits 0.033 under
the pathway backbone, without it 0.087 under, so the module recovers about two thirds of
what the pathway annotation supplies. The fold protocol (round 111, job 2699, one seed per
fold) reads 0.892 [0.83, 0.94] for the annotation-free backbone with the module against
0.795 [0.71, 0.88] without, beside 0.932 and 0.878 for the pathway backbone. The ten repeated splits (round 109) read the same
pair at three seeds per split; the read is fixed in the readout as the per-split seed-mean
difference and is reported in Table~\ref{tab:simplification-paired} as the cells complete;
a row of that table with fewer than ten splits or fewer than three seeds on every split is
partial.

\subsection{Rounds 84 to 98: what each simplification costs\secstatus{todo}}
\label{sec:round84-98}

Three simplifications make the module cheaper or easier to carry to another dataset:
reading one medium instead of four, so that a record stays one genotype in one
environment; dropping the stored flux vector, so that a record carries one growth ratio
per medium and no 4,131-column row; and dropping rule-count columns. A growth ratio is
itself one FBA program per genotype and medium, so dropping the vector saves the second
program of the parsimonious solve and the storage, not the solver. Rounds 84 to 98 (jobs 2616 to 2667, job 2622 aside) measure each against the
same baselines on the released split (Table~\ref{tab:simplification-ladder}) and, paired
by split, on the ten repeated outer splits of Section~\ref{sec:round75}
(Table~\ref{tab:simplification-paired}). The repeated-split arms that carry the solve were
extended to nine seeds per split (rounds 89 to 93 and 102 to 103), and the two recipes
with neither the vector nor four media to six (rounds 100 and 105). Round 101 swapped
the second medium and round 106 added the fold ensemble with no metabolism.

\paragraph{Every form of the module beats no metabolism.} Against the same cell with no
metabolic column, on the ten repeated splits at nine seeds each, the one-medium recipe
with the vector and the counts reads $+0.061$ and the four-media recipe $+0.059$, each
higher on ten of ten splits and on 89 and 88 of 90 split-and-seed pairs. The vector
without the counts reads $+0.040$, the counts without the vector $+0.041$ and the four
ratios alone $+0.028$, each on ten of ten. At six seeds the two-media counts recipe reads
$+0.036$ and the one-medium counts recipe $+0.032$, each on nine of ten. Round 106 gave
the no-metabolism cell its fold ensemble at three seeds per fold, 0.878, against 0.932
for the one-medium recipe under the same protocol.

\paragraph{One medium costs nothing on two protocols and about 0.01 on the third.} With
the flux vector and the counts kept, the YPD ratio alone against all four ratios reads
$+0.002$ over the ten splits at nine seeds each, higher on six of ten. The fold ensembles
agree, 0.932 against 0.934 at one seed per fold and 0.932 against 0.933 at three. The
cells of the released split do not: at round 83's six seeds they read 0.916 against
0.918, and round 107 (job 2683) took both recipes to nine seeds, where they read 0.912
against 0.920, $-0.009$ paired by seed with the one-medium cell higher on one of nine.
Two of the three added one-medium seeds read under 0.90. So round 83's conclusion holds
on ten other held-out fifths and as an ensemble, and on the released 195 genes the
single-seed cost is at most a hundredth.

\paragraph{The solve and the counts are complementary and about equal.} Over the four
ratios alone the counts are worth $+0.013$ (six of ten) and the vector $+0.012$ (eight
of ten) on the repeated splits, nine seeds each. Added to the other, the counts are worth
$+0.019$ (nine of ten) and the vector $+0.018$ (nine of ten). The released split orders them the other way,
0.911 for the vector over nine seeds against 0.902 for the counts over six, and so do the fold
ensembles, 0.920 against 0.909. Across the three protocols neither half is reliably the
larger, and the pair is reliably above either.

\paragraph{The simplifications do not stack.} Once the vector is gone the other media
carry signal the counts do not replace. With the counts and no flux vector, four ratios read
0.902 on the released split, the YPD and SC ratios 0.896, and the YPD ratio alone 0.888
(fold ensembles at three seeds per fold 0.909, 0.909 and 0.898),
and on the repeated splits at six seeds one medium for four costs $-0.011$ (higher on two
of ten) while two media for four cost $-0.006$ (two of ten), so the second medium buys
back about a third of the gap. Round 101 asked whether the second medium matters: the
YPD and SM ratios read 0.892 per cell and 0.907 as a one-seed fold ensemble, $-0.004$
against YPD and SC paired by seed (higher on two of six), so synthetic minimal medium
is no better a partner than synthetic complete. The recipe that respects one experiment and stores no flux vector,
one ratio with the six counts, reads 0.888 per cell and 0.902 as a fold ensemble against
0.920 and 0.934 for the full recipe, and $-0.029$ on the repeated splits, lower on ten of
ten. The cheapest admissible recipe that keeps the full score is therefore the
one-medium recipe with the vector. Its growth ratio already requires one linear program
per genotype and medium, run once offline, and the parsimonious solve that yields the
vector adds a second program on the same model, so the vector is nearly free once the
ratio is computed.

\paragraph{The six counts are redundant among themselves.} Round 77's misrank pattern
pointed at the sole-reaction count as the one useful column. Round 84 tested that: beside the vector, that count alone
reads 0.914, the other five without it 0.918, and all six 0.918, over six seeds each. The
five without it recover the whole gain ($-0.001$ against all six) and the one alone
recovers most of it ($+0.006$ over none, higher on five of six seeds), so no single
column carries it.

\paragraph{The training budget is not a limit.} Round 97 reran the full recipe and the
counts-alone recipe at 400 epochs. Validation selects the scored epoch, which falls
between 30 and 76 in all twelve cells, and each 400-epoch cell reproduced its 200-epoch
partner's trajectory, selected the same epoch and returned the same held-out AUROC to the
last digit. Nothing after epoch 76 is scored under this selection rule. That
reproducibility held within one code version on one card; ten fold cells of the
layer-free form rerun in round 92 under a later code version differ from their first runs
by up to 0.066 on a single cell, while the three-seed fold ensemble reads 0.921 against
0.919 on the first runs (round 34).

\paragraph{Three seeds per fold do not move the fold ensembles.} Rounds 92 and 94 took
each fold model from one seed to the mean logit of three. The ensembles read 0.933 for
the full recipe, 0.932 for one medium, 0.921 for the vector without counts and 0.909 for
the counts without the vector, each within 0.001 of its one-seed value.

%% notes-tex/028-gene-essentiality/tables/simplification_ladder.tex
%% GENERATED by experiments/028-gene-essentiality/scripts/notes_tex_tables.py
%% SOURCE: results/simplification_readout.json (ladder)
\begin{table}[H]
  \centering
  \footnotesize
  \caption[Rounds 83 to 98: media and solver on the released split]{What the affine columns cost and what they score, on the same backbone (pathway multi-hot with ProtT5). Media is the number of FBA growth ratios on the logit, each a different environment; solver says whether the recipe needs one linear program per genotype, run once offline. Columns: seeds, their mean held-out AUROC and its spread, and the five-fold ensemble at one seed per fold and at three. Best of each column in bold. Written by \texttt{notes\_tex\_tables.py} from \texttt{results/simplification\_readout.json}.}
  \label{tab:simplification-ladder}
  \begin{tabular}{@{}lrlrrrrr@{}}
    \hline
    affine columns & media & solver & seeds & cell mean & sd & folds, 1 seed & folds, 3 seeds \\
    \hline
    four ratios, flux vector, rule counts & 4 & yes & 6 & \textbf{0.918} & 0.006 & \textbf{0.934} & \textbf{0.933} \\
    YPD ratio, flux vector, rule counts & 1 & yes & 6 & 0.916 & 0.003 & 0.932 & 0.932 \\
    four ratios, flux vector & 4 & yes & 9 & 0.911 & 0.007 & 0.920 & 0.921 \\
    four ratios, rule counts & 4 & no & 6 & 0.902 & 0.005 & 0.909 & 0.909 \\
    YPD and SC ratios, rule counts & 2 & no & 6 & 0.896 & 0.013 & 0.905 & \multicolumn{1}{c}{--} \\
    YPD ratio, rule counts & 1 & no & 6 & 0.888 & 0.013 & 0.902 & \multicolumn{1}{c}{--} \\
    four ratios & 4 & no & 3 & 0.890 & 0.005 & \multicolumn{1}{c}{--} & \multicolumn{1}{c}{--} \\
    YPD ratio & 1 & no & 6 & 0.877 & 0.011 & \multicolumn{1}{c}{--} & \multicolumn{1}{c}{--} \\
    \hline
  \end{tabular}
\end{table}

%% notes-tex/028-gene-essentiality/tables/simplification_paired.tex
%% GENERATED by experiments/028-gene-essentiality/scripts/notes_tex_tables.py
%% SOURCE: results/simplification_readout.json (paired)
\begin{table}[H]
  \centering
  \footnotesize
  \caption[Rounds 84 to 99: the paired differences]{Every paired difference the simplification rounds fixed in advance, first recipe minus second. On the repeated splits the unit is a split, scored as the mean over the seeds both arms ran; the test blocks overlap, so no p-value is given. Columns: units, seeds per unit, the two means, the mean difference with its standard deviation over units, and the units on which the first recipe scored higher. Written by \texttt{notes\_tex\_tables.py} from \texttt{results/simplification\_readout.json}.}
  \label{tab:simplification-paired}
  \begin{tabular}{@{}p{66mm}rlrrrr@{}}
    \hline
    first minus second & units & seeds & first & second & difference (sd) & higher \\
    \hline
    \multicolumn{7}{@{}l}{\textit{ten repeated outer splits, unit = split}} \\
    one medium for four, solve and counts kept & 10 & 9 & 0.863 & 0.861 & $+0.0021$ (0.0071) & 6/10 \\
    counts for the flux vector, four media & 10 & 3 & 0.845 & 0.843 & $+0.0015$ (0.0246) & 5/10 \\
    counts over the four ratios alone & 10 & 3 & 0.845 & 0.832 & $+0.0122$ (0.0204) & 7/10 \\
    flux vector over the four ratios alone & 10 & 3 & 0.843 & 0.832 & $+0.0107$ (0.0233) & 7/10 \\
    counts added to the flux vector, four media & 10 & 3 & 0.862 & 0.843 & $+0.0192$ (0.0198) & 8/10 \\
    no solve, counts kept, four media & 10 & 6 to 9 & 0.844 & 0.861 & $-0.0175$ (0.0154) & 1/10 \\
    one medium for four, no solve & 9 & 3 & 0.832 & 0.843 & $-0.0114$ (0.0123) & 2/9 \\
    no solve and one medium, against the full recipe & 9 & 3 & 0.832 & 0.862 & $-0.0297$ (0.0162) & 0/9 \\
    \hline
    \multicolumn{7}{@{}l}{\textit{rule-count subsets on the released split, unit = seed}} \\
    sole-reaction count alone, against all six & 6 & 1 & 0.914 & 0.918 & $-0.0046$ (0.0088) & 2/6 \\
    five without it, against all six & 6 & 1 & 0.917 & 0.918 & $-0.0010$ (0.0076) & 3/6 \\
    sole-reaction count alone, against none & 6 & 1 & 0.914 & 0.908 & $+0.0056$ (0.0114) & 5/6 \\
    five without it, against none & 6 & 1 & 0.917 & 0.908 & $+0.0092$ (0.0068) & 5/6 \\
    \hline
    \multicolumn{7}{@{}l}{\textit{400 epochs against 200 on the released split, unit = seed}} \\
    flux vector and counts & 6 & 1 & 0.918 & 0.918 & $+0.0000$ (0.0000) & 0/6 \\
    counts alone & 6 & 1 & 0.902 & 0.902 & $+0.0000$ (0.0000) & 0/6 \\
    \hline
  \end{tabular}
\end{table}

